% Section 5: the venue panel, the report's original dataset, plus the
% quasi-experimental tests. The key analytic move is separating the aggregate
% decline (real) from per-venue performance (no different from other bars).
\section{Austin's live-music capacity is shrinking through room loss rather than weaker trade}
\label{sec:venues}

Musicians need places to play, and Austin's rooms are measurable in a way
musicians themselves are not. The Texas Comptroller of Public Accounts collects
a monthly filing from every business in the state holding a mixed-beverage
permit (the state's license to sell liquor by the drink), reports each
location's gross receipts, and publishes the file on the state open data
portal.\footnote{Texas Comptroller of Public Accounts, \textit{Mixed Beverage
Gross Receipts}, Texas Open Data Portal dataset \texttt{naix-2893}, at
\url{https://data.texas.gov/dataset/Mixed-Beverage-Gross-Receipts/naix-2893},
retrieved 2 August 2026. The population is every permittee filing under the
mixed-beverage permit categories (MB, BQ, PE and their subtypes); the file is
refreshed monthly and the extract used here covers obligation end dates January
2007 through August 2026. The dataset omits rooms operating on a beer-and-wine
permit only, and it omits venues without a separable alcohol permit or folded
into a hotel or nonprofit permit, so it undercounts several well-known Austin
listening rooms; the panel construction below inherits those omissions. Every
filing is a legal tax record rather than a survey response, which makes it a
census of the licensed universe rather than a sample of it, and each row carries
the permittee's taxpayer name, the location address and the reported gross
receipts for the month.} This study matches those filings to named venues by
taxpayer name and location address and follows a panel of
\NVenuePanelIsCurated\ Austin live-music rooms monthly from January 2007
through June 2026, tracked across changes of owner and legal name.\footnote{The panel is a
\emph{curated} tracking set of named Austin live-music rooms, not a census of
the licensed universe and not a probability sample of it. Rooms enter by name
identification against a research list of Austin live-music venues, so the panel
over-represents established rooms in the central city (Red River, downtown,
East Austin and South Congress), and it excludes venues without a separable
liquor permit, which removes several well-known listening rooms including
Cactus Cafe, Bass Concert Hall, Central Presbyterian, Carousel Lounge,
Meanwhile Brewing, Whip In and Geraldine's. Receipts also measure alcohol
sales rather than music activity: a room's mixed-beverage receipts say nothing
directly about what performers were paid, only what the bar took at the register
that month. Ten festival concession venue-months, in which a festival company
filed under a venue-linked entity name for a month at a scale more than three
hundred times the host venue's median month, are excluded from every venue-level
figure that follows and are worth \NVenueFestivalDollarsExcluded\ in nominal
dollars; the citywide totals used for comparison keep them, because there the
money is real Austin licensed-beverage activity and only its venue label is
misleading. The panel is compared against the full Travis County and City of
Austin licensed-beverage universe throughout this section, and each comparison
names the two sides in matching terms.} With that panel this section asks
whether the sector is contracting because surviving rooms trade worse than the
rest of the licensed-beverage market, or because rooms close and no one replaces
them.

\subsection{Receipts fell at live-music rooms while the market around them added locations}

\begin{figure}[htbp]
\centering
\includegraphics[width=\textwidth]{fig09_venue_receipts_index.png}
\caption{Real mixed-beverage receipts indexed to 2019, and the count of active
permit locations, Travis County and City of Austin, 2013--2025, 2025 dollars.
Source: Texas Comptroller of Public Accounts, \textit{Mixed Beverage Gross
Receipts}, Texas Open Data Portal dataset \texttt{naix-2893}, matched to named
venues by taxpayer name and location address and deflated with the national
Consumer Price Index for All Urban Consumers (CPI-U, series
\texttt{CUUR0000SA0}). How to read it: the upper panel compares three groups
against their own 2019 level, so all three start at 100 by construction; the
lower panel shows how many distinct permit locations filed receipts in each
year. The three groups are the \NVenuePanelIsCurated-venue curated live-music
panel; its core live-music tier of rooms whose main business is live
performance; and the full Travis County and City of Austin mixed-beverage
universe. Caveat: receipts measure alcohol sales, not attendance, ticket
revenue or what performers were paid; the curated panel is not a census, so it
is smaller and more central than the universe by construction.}
\label{fig:venueindex}
\end{figure}

Real receipts fell between 2019 and 2025 at the rooms the curated panel
tracks: \NVenuePanelPctTwoZeroOneNineTwoZeroTwoFive\ across the
\NVenuePanelIsCurated-room panel and
\NVenueCorePctTwoZeroOneNineTwoZeroTwoFive\ in its core live-music tier, the
rooms whose main business is live performance.\footnote{Every filing is
deflated before anything is added up:
\(R^{\mathrm{real}}_{it} = R_{it} \times d_{y(t)}\), where \(R_{it}\) is venue
\(i\)'s nominal receipts in month \(t\), \(y(t)\) is the calendar year
containing that month, and \(d_y = C_{2025}/C_y\) is the ratio of national
consumer price index annual averages. Group totals are
\(S_{g,y} = \sum_{i \in g} \sum_{t \in y} R^{\mathrm{real}}_{it}\), and the
change quoted is \(100\,(S_{g,2025}/S_{g,2019} - 1)\).
Figure~\ref{fig:venueindex} plots the same totals rebased on 2019,
\(I_{g,y} = 100 \times S_{g,y}/S_{g,2019}\), which is why all three lines
start at 100. These are totals over an unbalanced set of filers, so a group
can fall because rooms closed as well as because open rooms took less; that
distinction is the subject of the next subsection. Restricted to the
\NVenueBalancedN\ venues filing in both years, the change is
\NVenueBalancedPctTwoZeroOneNineTwoZeroTwoFive, and
\NVenueShrankBalancedN\ of them were smaller in real terms in 2025.} The full
Travis County and City of Austin mixed-beverage universe moved in the opposite
direction over the same six years, changing
\NVenueCitywidePctTwoZeroOneNineTwoZeroTwoFive\ in real receipts, and the count
of distinct permit locations filing during the year grew from
\NVenueCitywidePermitsTwoZeroOneNine\ in 2019 to
\NVenueCitywidePermitsTwoZeroTwoFive\ in 2025.\footnote{The two sides of that comparison
are not built alike. The ten festival-aggregation venue-months are stripped out
of every venue-level series but stay in the citywide market total, where the
money is genuine licensed-beverage activity and only its venue label is wrong.
On the same ex-festival basis the citywide change is about +1.6 percent rather
than \NVenueCitywidePctTwoZeroOneNineTwoZeroTwoFive. The citywide total also
absorbs \NVenueCitywidePermitsPctTwoZeroOneNineTwoZeroTwoFive\ more permit
locations over the period, so it is not a per-room comparison; on a common
basis, the tracked panel held \NVenuePanelShareTwoZeroOneNine\ of citywide
receipts in 2019 and \NVenuePanelShareTwoZeroTwoFive\ in 2025.} Austin kept
adding places to drink while its dedicated live-music rooms appeared to
contract. Some live performance may also have shifted toward informal or
non-dedicated settings this alcohol-tax panel cannot observe---coffee shops,
restaurants, and breweries that host occasional acts---which typically pay
performers less and less regularly, and are not where full-time musicians earn
most of their income.

A separate City of Austin administrative record moves the same
direction.\footnote{City of Austin, \textit{Sound Ordinance Permits}, Socrata
dataset \texttt{ryu3-tuin} at
\url{https://data.austintexas.gov/resource/ryu3-tuin.csv}, retrieved 1 August
2026. The Music \& Entertainment Division issues an Outdoor Music Venue (OMV)
permit under City Code Chapter 9-2 to any establishment that presents amplified
sound outdoors, and the by-year counts used here are annual issuances
tabulated from that permit file. The series is independent of the Comptroller
receipts data because it is a city licensing series rather than a state tax
filing, and it moves in the same direction as the receipts panel; that
concurrence is why it is reported here.} The Music \& Entertainment Division
issued \NVenueOmvPeak\ Outdoor Music Venue permits in both 2014 and 2017 at
the series peak and \NVenueOmvTwoZeroTwoFive\ in 2025, a change of
\NVenueOmvPctFromPeak\ from the peak.\footnote{The change is
\(100\,(P_{2025}/P_{\mathrm{peak}} - 1)\), where \(P_y\) is the count of
outdoor music venue permits the City issued in calendar year \(y\). A permit
count is not a dollar figure, so nothing is deflated, and one room can hold
several permits across years, so the series measures permitted outdoor music
activity rather than a count of distinct rooms.}

\begin{figure}[htbp]
\centering
\includegraphics[width=\textwidth]{fig11_venue_supply.png}
\caption{Outdoor Music Venue permits issued by the City of Austin each
calendar year, with net venue entries minus exits in the curated live-music
panel. Sources: City of Austin, \textit{Sound Ordinance Permits}, Socrata
dataset \texttt{ryu3-tuin}; Texas Comptroller of Public Accounts, \textit{Mixed
Beverage Gross Receipts}, dataset \texttt{naix-2893}. How to read it: bars show
permits issued on the left axis and the line shows net additions (entries
minus exits) to the curated \NVenuePanelIsCurated-room panel on the right axis;
the 2026 bar covers only the partial year available at extract. Caveat: a
permit is not a venue, because one room may hold several across years, so the
series measures permitted outdoor music activity rather than a count of rooms.}
\label{fig:venuesupply}
\end{figure}

\subsection{Surviving rooms lost ground, but no more than the rest of the market did}

The aggregate gap invites an obvious explanation, that live-music rooms are
weaker businesses than other bars. Rooms that stayed open did trade down:
among the \NVenueBalancedN\ venues filing in both 2019 and 2025, real receipts
changed \NVenueBalancedPctTwoZeroOneNineTwoZeroTwoFive\ and
\NVenueShrankBalancedN\ shrank. The 2019-to-2025 change by itself cannot tell
whether the surviving live-music rooms fell further than the bars around them
did, because the same 2019-to-2025 change at a citywide comparison group would
carry the same seasonal weather, market-wide inflation and pandemic exposures
that also hit the live-music rooms. The divergence test below strips those
common shocks out, by holding each venue and each calendar month constant, and
asks what is left after both sides have been controlled against themselves.

Comparing the core live-music tier against the rest of the Travis County and
City of Austin mixed-beverage market before and after March 2020, with each
venue and each month held constant, gives an average difference of
\NVenueDidPct, and the 95 percent interval runs from about 14 percent below
the rest of the market to about 17 percent above it (coefficient
\NVenueDidCoef\ log points, clustered standard error \NVenueDidSe, $p =$
\NVenueDidP), across \NVenueDidN.\footnote{The specification is a two-way fixed-effects regression, meaning a
linear regression that includes one intercept for every venue and one
intercept for every calendar month, which together absorb any fixed
characteristic of a room and any month common to the whole market. The
dependent variable is log real monthly receipts, the regressor of interest is
an indicator for core live-music venues interacted with the post-March-2020
period, and standard errors are clustered by venue (that is, they allow the
residuals from the same room to be correlated across months, and treat venues
as the independent units):
\[
\ln R^{\mathrm{real}}_{it}
  = \alpha_i + \tau_t
  + \beta\,\bigl(\mathrm{Core}_i \times \mathrm{Post}_t\bigr)
  + \varepsilon_{it},
\]
where \(R^{\mathrm{real}}_{it}\) is venue \(i\)'s mixed-beverage receipts in
month \(t\) in 2025 dollars; \(\mathrm{Core}_i\) is one for a core live-music
room and zero for other mixed-beverage permit locations in Travis County or the
City of Austin; the panel's live-music-adjacent tier is dropped from the
estimation sample rather than coded zero, because those rooms programme live
music intermittently and belong to neither the treated group nor the
comparison group; \(\mathrm{Post}_t\) is one from March 2020 onward; \(\alpha_i\) is a
venue fixed effect that absorbs everything constant about that room including
its size, address and tier; \(\tau_t\) is a year-month fixed effect that
absorbs everything common to the whole market in that month including the
pandemic, market-wide inflation and season; and \(\varepsilon_{it}\) is the
residual. Only the interaction is estimable, because \(\alpha_i\) absorbs
\(\mathrm{Core}_i\) and \(\tau_t\) absorbs \(\mathrm{Post}_t\). The
interaction is built as an explicit product variable rather than entered
factorially: entered factorially it is collinear with the absorbed effects and
Stata's \texttt{reghdfe} drops the term of interest. Estimation covers
\NVenueDidN\ from January 2013 to December 2025. Months with zero receipts
leave a log specification, so the coefficient describes venues conditional on
trading at all. It converts to a proportional difference as
\(100\,(e^{\hat\beta} - 1)\). This is a descriptive divergence test: no policy
switches on in March 2020 for one group and not the other, and nothing here is
randomised. It is also a precise null rather than an absence of evidence: the
95 percent interval runs from about $-13.7$ to $+16.5$ percent, so a
difference between core live-music rooms and the rest of the mixed-beverage
market much larger than that in either direction is not consistent with these
data.} An event study, which is the same regression with a separate coefficient for
each calendar year rather than one averaged post-March-2020 coefficient, gives
the same reading year by year. Core live-music rooms dropped about 54 log
points below the comparison group in 2020, returned to within statistical
distance of it by 2022, and have tracked it since. The 2019 base is not a
neutral reference, because the pre-2019 coefficients drift upward rather than
sitting flat, so a counterfactual that continued that drift would sit above
zero after 2022 rather than at it.\footnote{The event study replaces the
single interaction with one term per calendar year:
\[
\ln R^{\mathrm{real}}_{it}
  = \alpha_i + \tau_t
  + \sum_{y \neq 2019} \beta_y\,
    \bigl(\mathbf{1}\{\mathrm{year}(t) = y\} \times \mathrm{Core}_i\bigr)
  + \varepsilon_{it},
\]
where \(\mathbf{1}\{\mathrm{year}(t) = y\}\) is one when month \(t\) falls in
calendar year \(y\), 2019 is omitted, and every other symbol, the estimation
sample and the clustering are as in the previous footnote. Each \(\beta_y\)
therefore reads as the gap between core live-music rooms and the rest of the
market in year \(y\), measured against their own 2019 gap, and the 2020 term is
the one quoted above. The year terms are constructed as explicit products with
2019 dropped by hand: asking the estimator to hold 2019 as the base is not
honoured once venue and year-month effects absorb the alternatives, and it
silently renormalises the whole path onto a different year. Pre-2019
coefficients drift upward instead of sitting flat, so parallel trends do not
hold cleanly and the design is descriptive, not causal. A fourth specification
in the exported table replaces the log with the inverse hyperbolic sine, which
keeps zero-receipt months in the sample, and it returns a large negative
coefficient rather than a null; that column answers a different question,
conditioning on nothing rather than on trading, and no number quoted here comes
from it, but a reader should know it points the other way. The full
coefficient path is in the technical appendix.}

If the rooms that stay open trade no worse than other bars, the rest of the
aggregate decline has to come from which rooms exist. \NVenueExitsTotal\ of the
\NVenuePanelIsCurated\ tracked venues had stopped filing altogether by June
2026. \NVenueExitsTwoZeroTwoZeroTwoZeroTwoSix\ of those last filings fall in
2020 or later, and \NVenueExitsTwoZeroTwoZeroTwoZeroTwoSixCore\ of the post-2020
exits were core live-music rooms. Net additions to the panel (entries minus
exits) came to \NVenueNetAddsTwoZeroZeroSevenTwoZeroOneFive\ over 2007--2015 and
\NVenueNetAddsTwoZeroOneSixTwoZeroTwoSix\ over 2016--2026.\footnote{The two
net-add figures are not built alike. The Comptroller receipts file begins in
January 2007, so the 44 venues already trading then all register as 2007
entries; excluding 2007 entirely, the comparable 2008-to-2015 net is +36. The
later figure is not left-censored but is right-truncated at June 2026. Ceasing
to file usually means closure but can also mean a permit transfer, an entity
restructuring, or a switch to a beer-and-wine permit, so these counts are
stopped reporting rather than verified closures.}

\begin{keyfinding}
{\small\textbf{This distinction matters for policy.} Austin's live-music
capacity is contracting because rooms stop filing and are not replaced, while
the wider licensed-beverage market keeps adding permit locations. On
receipts alone, rooms that stayed open trade at the same rate against 2019
as comparable non-music bars; the gap that separates the panel from the
citywide total is not a gap in per-room revenue among survivors. The
receipts file records what bars sell, not what venues spend, so it cannot
say whether operating margins hold; it says only that the aggregate decline
tracks which rooms exist rather than how the open ones sold liquor.}
\end{keyfinding}

Panel venues with lower real receipts in 2025 than in 2019 appear in every
part of Austin the panel can name (Figure~\ref{fig:venuedecline}). The ten
largest dollar declines are concentrated where the panel's own coverage is
concentrated: six sit downtown, on Rainey Street or in the Red River district,
and the other four are in South Austin, in South Central and near the
University of Texas campus. Because dollar size tracks venue size and the
curated panel skews central by construction, that list of ten largest declines
cannot be read as evidence about where decline is or is not spread within the
city.

\begin{figure}[htbp]
\centering
\includegraphics[width=\textwidth]{fig10_venue_map.png}
\caption{Venues in the curated \NVenuePanelIsCurated-room panel with lower
real receipts in 2025 than in 2019, by part of Austin, and the ten largest
dollar declines by venue name. Source: Texas Comptroller of Public Accounts,
\textit{Mixed Beverage Gross Receipts}, dataset \texttt{naix-2893}, deflated
with the national CPI-U. How to read it: each row on the left gives the number
of panel venues that shrank against the number that filed in 2019, so the
denominators differ by area and several rows rest on only three venues; the
right panel shows the size of the ten largest dollar declines by venue name.
Caveats: named parts of Austin are assigned from venue addresses and are
descriptive groupings rather than official boundaries; panel venues that
stopped filing count as zero in 2025 and so appear as declines, which covers 16
of the 64 counted here, and three of the ten largest named declines are rooms
with no 2025 receipts at all.}
\label{fig:venuedecline}
\end{figure}

\begin{figure}[htbp]
\centering
\includegraphics[width=0.82\textwidth]{fig10b_district_map.png}
\caption{Outdoor Music Venue sound permits issued by the City of Austin,
2009--2018, by council district. Source: City of Austin, \textit{Sound
Ordinance Permits, geographic extract}, Socrata dataset \texttt{g3rj-dfgm} at
\url{https://data.austintexas.gov/resource/g3rj-dfgm.csv}, retrieved 1 August
2026; council district boundaries from City of Austin open data. How to read
it: darker districts issued more permits, and the class breaks are set by hand
because the permit distribution is too concentrated in District 9 for equal
intervals to separate anything. Caveats: annual renewals count as separate
permits, so one room can appear many times and the map shows permitted
activity, not distinct rooms; permits are assigned to the council district
boundaries now in force, which took effect in January 2015, so the 2009-to-2014
permits are mapped retrospectively onto districts that did not exist when they
were issued; and the series stops at 2018 because the geocoded file contains
only a partial extract for 2019 (13 rows against the 99 the City by-year count
reports).}
\label{fig:districtmap}
\end{figure}

Between 2009 and 2018, permitted outdoor music activity concentrated in
Council District 9, which covers downtown Austin and the central corridors,
and District 9 accounted for \NVenueOmvPermitsDistrictNineShare\ of the
\NVenueOmvPermitsGeoTwoZeroZeroNineTwoZeroOneEight\ Outdoor Music Venue permits
the City issued in those years. Citywide, the Music \& Entertainment Division's
Outdoor Music Venue permits have since fallen \NVenueOmvPctFromPeak\ from
their 2014--2017 peak, so the geography of permitted outdoor music activity may
have moved since 2018. The curated receipts panel cannot test that
distributional question independently, because its own coverage skews central
by construction.

A separate news-derived timeline of \NVenueClosuresDatedN\ Austin
live-music venue closures reported between January 2015 and June 2026, coded
by primary reason, puts property pressure level with the pandemic:
\NVenueClosuresCovid\ closures cite COVID and
\NVenueClosuresProperty\ cite a property cause, meaning a rent increase, a
sale of the building, a lease expiry or redevelopment.\footnote{The timeline
lists \NVenueClosuresDatedN\ Austin live-music venue closures dated between
January 2015 and June 2026, coded by primary reason from local news coverage
in \emph{Austin Chronicle}, \emph{Austin American-Statesman}, KUT/KUTX,
\emph{Austin Monitor} and CultureMap Austin, plus follow-up items from
\emph{Do512} and \emph{Texas Monthly}, compiled for this study and released
as the project evidence file
\texttt{01\_evidence/08\_venues\_ecosystem/venue\_timeline.csv}. Reasons are
journalist-cited rather than audited, only the primary reason is coded when a
venue reports more than one, and \NVenueClosuresNoReason\ closures have no
reason stated in the sources examined. The timeline is a news-derived record
rather than a register, so it misses venues the research team did not find
coverage for, and an absence here means not found rather than nothing happened.
A separate cross-check against the Comptroller mixed-beverage file found
\NVenueClosureCrossvalContinue\ venues still filing receipts after their
reported closure date, which usually means the programming ended before the
business did.} Respondents to the 2022 Greater Austin Music Census who present
live music, most of whom are not primary venue operators, ranked property tax
first among the nine business pressures the census asked them to rank
(Section~\ref{sec:census}). The two records both point at the cost and security
of the building, though they measure different things: a tax on owners in the
census answer, and rent, sale and tenure in the news timeline.

\subsection{Three tests find no detectable change in the receipts}

Each test below compares the panel venues exposed to a particular intervention
against the rest of the mixed-beverage panel. The panel is the natural setting
for these tests because each intervention hit a small, named set of Austin
live-music rooms and left a monthly, room-level receipts trail on either side
of the event, which is exactly what a before-and-after comparison against a
non-treated group needs.

\textbf{The 2017 extended-hours pilot.} Five named Red River venues ran later
outdoor music hours in 2017 under a City of Austin sound-ordinance pilot; the
only ordinance text located names the pilot as a six-month experiment later
extended to 30 April 2018.\footnote{City of Austin, Law Department, \textit{Draft
ordinance amending City Code Chapter 9-2 (Red River extended-hours pilot)},
dated 30 March 2018, at
\url{https://services.austintexas.gov/edims/document.cfm?id=296553}, retrieved
1 August 2026. The City evaluation memo referenced in the press has not been
located, so the pilot window used in the reported city evaluation (May to
November 2017) and the window the Law Department draft ordinance describes do
not match; both are reported below.} Press accounts of a City evaluation
reported a \NQuasiRedriverClaimUnverified\ rise in payments to local musicians,
alongside 10 percent more local and regional acts booked and 13 percent fewer
total tickets sold; the evaluation memo behind those figures has never been
located and nobody has independently checked them. The Comptroller receipts
contain no matching change in venue revenue.
Once venue-specific seasonality is absorbed, which matters because six of the
pilot window's seven months fall in the second half of the year, the estimate
is \NQuasiRedriverDidSeasonadj\ log points. The unadjusted specification gives
\NQuasiRedriverDidCoef\ log points but fails its own pre-trend test.
Randomization inference on the unadjusted specification, which asks how large
an estimate the same regression produces when five other panel venues are
labelled as treated at random, gives an exact $p =$
\NQuasiRedriverRiP\ across 5,000 placebo draws.\footnote{The specification is
\(\ln R^{\mathrm{real}}_{it} = \alpha_i + \tau_t
   + \beta\,(\mathrm{Pilot}_i \times W_t) + \varepsilon_{it}\),
where \(\mathrm{Pilot}_i\) is one for the five named venues, \(W_t\) is one
for the months of the reported pilot window from May to November 2017, and
\(\alpha_i\) and \(\tau_t\) are venue and year-month fixed effects. The sample
is the panel venues from January 2015 to December 2019 with the large ticketed
rooms excluded. With five treated venues a cluster-robust standard error is
not trustworthy, so inference is by randomization instead: the treatment dates
are held fixed, five venues are drawn at random without replacement from the
panel venues filing in nearly every month of the window, the same regression is
refitted on the placebo assignment, and the exact p-value is
\[
p = \frac{1 + \#\{\,r \le R : |\hat\beta_r| \ge |\hat\beta|\,\}}{R + 1},
\]
where \(\hat\beta\) is the estimate under the assignment that actually
happened, \(\hat\beta_r\) is the estimate under placebo assignment \(r\), and
\(R\) is 5,000, the number of draws. Adding one to numerator and denominator
counts the observed assignment among the possibilities, which is what keeps
the test exact rather than approximate. The bound quoted next comes from that
same placebo distribution: 90 percent of placebo estimates fall inside roughly
plus or minus 14 percent, so an effect larger than that would have shown. The
five venues were named into the pilot rather than drawn into it, so this is
descriptive.} Nine in ten placebo assignments produce estimates inside roughly
plus or minus 14 percent, so a receipts effect much larger than that would
probably have been visible here.\footnote{Pre-trends fail as first specified (p =
\NQuasiRedriverPretrendP), but the cause is seasonality rather than divergence:
the pilot window was six-sevenths second-half months. The pre-trend check
replaces the single interaction with half-year terms,
\(\ln R^{\mathrm{real}}_{it} = \alpha_i + \tau_t
   + \sum_{h \neq h_0} \pi_h\,(\mathbf{1}\{\mathrm{halfyear}(t) = h\}
   \times \mathrm{Pilot}_i) + \varepsilon_{it}\),
with the second half of 2016 as the omitted base \(h_0\), and tests jointly
whether the three half-years before the pilot are all zero,
\(H_0: \pi_{2015\mathrm{H}1} = \pi_{2015\mathrm{H}2}
   = \pi_{2016\mathrm{H}1} = 0\),
by a Wald test with standard errors clustered by venue. Absorbing
venue-specific seasonality, by adding a venue-by-month-of-year fixed effect
that lets each room have its own December and its own March rather than only
the seasonality the whole market shares, clears the pre-trend (p =
\NQuasiRedriverPretrendPSeasonadj) and moves the estimate to
\NQuasiRedriverDidSeasonadj. The reported May-to-November
window also conflicts with the only ordinance text in hand, which describes a
six-month pilot later extended to 30 April 2018.} \textbf{This does not disprove
the \NQuasiRedriverClaimUnverified\ claim.} Mixed-beverage receipts measure
alcohol sales, not the
door split between venue and performer, so a venue could raise what it pays
performers by a fifth without changing its bar takings.

\textbf{The 2024 Live Music Fund venue grants.} The City of Austin
awarded \NQuasiLmfTwoFourVenueAwards\ Live Music Fund venue grants in fiscal
2024, all at the same dollar amount.\footnote{City of Austin, \textit{Austin
Live Music Fund, FY2024 Awardee List}, at
\url{https://austin.widen.net/content/bpzuacfkk8/pdf/FY24-AwardeeList_Austin-Live-Music-Fund.pdf},
retrieved 1 August 2026. All \NQuasiLmfTwoFourVenueAwards\ awards were made at
\$60,000. The full FY2024 award total, adding musician and non-venue promoter
awardees, is 136 awards / \$4,365,000 by independent line-item tally. Program
mechanics and the fund's revenue source are the subject of
Section~\ref{sec:city}.} Against other tracked panel venues the estimate is
\NQuasiLmfTwoFourPostCoef\ log points, about 9 percent, with an exact
randomization-inference $p =$ \NQuasiLmfTwoFourRiP\ across 2,000 placebo draws.
No divergence between grantees and non-grantees is detectable in the years
before the awards ($p =$ \NQuasiLmfTwoFourPretrendP), though that window spans
the pandemic and the parallel-paths test has little power with 15 treated
venues; restricted to 2018 and 2019, the same joint test gives $p =$
\NQuasiLmfTwoFourPretrendPPrepandemic.\footnote{The specification is
\(\ln R^{\mathrm{real}}_{it} = \alpha_i + \tau_t
   + \beta\,(\mathrm{Grant}_i \times \mathrm{Post2024}_t) + \varepsilon_{it}\),
where \(\mathrm{Grant}_i\) is one for a fiscal 2024 venue grantee and
\(\mathrm{Post2024}_t\) is one from January 2024 onward, on monthly
observations from January 2018 to December 2025 with the large ticketed rooms
excluded and standard errors clustered by venue. Grantees with too few filings
in the base year to have a pre-period, and the one large room among them,
leave the estimation sample, so the regression rests on fewer venues than
received awards. The parallel-paths check is an annual event study with 2023
as the omitted base and a joint Wald test that the pre-award year terms are
all zero, \(H_0: \beta_{2018} = \beta_{2019} = \beta_{2020} = \beta_{2021} =
\beta_{2022} = 0\); that window spans the pandemic, so it asks whether the two
groups followed the same path through it; the narrower 2018-and-2019 version
quoted in the text rests on the 12 of 15 grantee venues already filing in 2019.
Randomization inference draws
the same number of placebo grantees from the venues meeting the same
filing-coverage requirement and computes the exact p-value by the formula
above, over 2,000 draws.} Grants are competitively scored, so recipients
differ from non-recipients in whatever the scoring rewarded, and the estimate is
an association rather than a program effect. The
\NQuasiLmfFyTwoSixNotEvaluable\ of the fiscal 2026 round cannot
be tested at all: the city made those awards in March 2026 and the receipts
panel ends that June.

\textbf{The Moody Center.} The Moody Center, a 15,000-seat multipurpose arena
on the University of Texas at Austin campus, opened for public events in April
2022.\footnote{Moody Center, \textit{About}, at
\url{https://www.moodycenteratx.com/about}, retrieved 1 August 2026; opening
event reported in local news coverage the same month. The arena competes for
concert bookings at the top of the market rather than at the size of the
tracked live-music rooms, so any spillover to the venue panel would appear as
a proximity gradient in receipts, not as a shift in the market average.} The
test asks whether panel venues near the arena traded differently from panel
venues farther out after it opened. Within four miles of the arena the
post-opening pattern does not vary with distance ($p =$
\NQuasiMoodyJointP). The design cannot say more than that: the comparison
group is only six venues more than four miles out, every distance-band standard
error is large, and no band is individually distinguishable from zero.\footnote{All three bands enter one regression,
\[
\ln R^{\mathrm{real}}_{it}
  = \alpha_i + \tau_t
  + \sum_{b=1}^{3} \beta_b\,
    \bigl(\mathrm{Band}_{bi} \times \mathrm{Post}_t\bigr)
  + \varepsilon_{it},
\]
where \(\mathrm{Band}_{bi}\) is one when venue \(i\) sits under one mile, one
to two miles, or two to four miles from the arena, venues more than four miles
away are the omitted reference group, and \(\mathrm{Post}_t\) is one from
April 2022 onward. Distance is straight-line miles from the arena's
coordinates, computed by the great-circle formula
\(D = 2\rho\arcsin\sqrt{\sin^{2}(\Delta\phi/2)
   + \cos\phi_1\cos\phi_2\sin^{2}(\Delta\lambda/2)}\),
with \(\phi\) latitude, \(\lambda\) longitude, and \(\rho\) the Earth's radius
in miles. The pre-period is January 2018 to December 2019 and the post-period
July 2022 to December 2025, so neither period includes the pandemic. The quoted
p-value is a Wald test of \(H_0: \beta_1 = \beta_2 = \beta_3\), which asks
whether the post-opening pattern varies with distance at all rather than
whether any single band moved; a continuous version replacing the bands with
post-opening times log distance is registered separately, so nothing rests on
where the band cuts fall. Standard errors are clustered by venue. Distance is
not assigned at random: near venues are central-city rooms and differ from
suburban ones in many ways at once, so this is a description of where the
divergence sits. Note also that any market-wide arena effect is absorbed by the
year-month fixed effects by construction, so this design tests proximity and not
the arena's overall effect on venue receipts.}

These receipts describe what rooms sold at the register. They do not describe
who played in them. Section~\ref{sec:census} turns to the only current
Austin-specific measurement of the people who work in those rooms, the 2022
Greater Austin Music Census, and takes up the retention question the receipts
panel cannot answer. Federal pandemic relief, the one intervention with a
visible signature in the mixed-beverage receipts, waits until
Section~\ref{sec:state}, because that signature traces the eligibility rule
rather than the program's effect and belongs with the state and national
programs.
